% Luc Destombe --- licence CC BY-NC-SA 4.0
% https://creativecommons.org/licenses/by-nc-sa/4.0/deed.fr
% Fichier autonome : compiler avec pdflatex, deux fois (signets, nombre de pages).
\documentclass[a4paper,10pt]{article}
\makeatletter
\newif\iflycee@niveaux
\lycee@niveauxtrue
\newif\iflycee@palatino
\lycee@palatinotrue

\RequirePackage[utf8]{inputenc}
\RequirePackage[T1]{fontenc}
\RequirePackage[french]{babel}
\RequirePackage{etoolbox}

\newcommand{\ShorthandsOff}{\@ifundefined{shorthandoff}{}{\shorthandoff{;:!?}}}
\newcommand{\ShorthandsOn}{\@ifundefined{shorthandon}{}{\shorthandon{;:!?}}}
\ShorthandsOff
\AtBeginDocument{\ShorthandsOn}
\AtBeginEnvironment{tikzpicture}{\ShorthandsOff}

\RequirePackage{geometry}
\RequirePackage{fancyhdr}
\RequirePackage{lastpage}
\RequirePackage{multicol}
\RequirePackage{array}
\RequirePackage{enumitem}
\RequirePackage{xcolor}
\RequirePackage{needspace}

\RequirePackage{amsmath,amssymb}
\RequirePackage{mathpazo}
\DeclareMathSizes{10.95}{10.95}{8}{5.5}
\begingroup\catcode`\_=\active \catcode`\^=\active
\gdef_#1{\sb{\scriptscriptstyle #1}}
\gdef^#1{\sp{\scriptscriptstyle\lycee@moinsepais #1}}
\endgroup
\newcommand\lycee@moins{\mathbin{%
  \setbox\z@\hbox{$\m@th\scriptscriptstyle\lycee@moinsorig$}%
  \dimen@=.55\wd\z@ \advance\dimen@-.5pt
  \hbox to.75\wd\z@{\hss$\m@th\scriptscriptstyle\vcenter{\hbox{%
    \tikz\draw[line width=.5pt,line cap=round](0,0)--(\the\dimen@,0);}}$\hss}}}
\begingroup\lccode`\~=`\- \lowercase{\endgroup\let~}\lycee@moins
\newcommand\lycee@moinsepais{\mathcode`\-="8000 }
\AtBeginDocument{\mathchardef\lycee@moinsorig=\mathcode`\-\relax}
\AtBeginDocument{\catcode`\_=12 \mathcode`\_="8000
  \catcode`\^=12 \mathcode`\^="8000 }
\newcommand\lycee@exposants{%
  \fontdimen13\textfont2=.48\fontdimen6\textfont2
  \fontdimen14\textfont2=.48\fontdimen6\textfont2
  \fontdimen15\textfont2=.48\fontdimen6\textfont2
  \fontdimen13\scriptfont2=.48\fontdimen6\scriptfont2
  \fontdimen14\scriptfont2=.48\fontdimen6\scriptfont2
  \fontdimen15\scriptfont2=.48\fontdimen6\scriptfont2}
\every@math@size\expandafter{\the\every@math@size\lycee@exposants}
\newlength{\retraitcalculs}
\setlength{\retraitcalculs}{2em}

\RequirePackage{tikz}
\usetikzlibrary{arrows.meta,calc,babel,shapes.misc}
\RequirePackage[most]{tcolorbox}
\tcbuselibrary{skins,breakable}

\definecolor{coulDef}{HTML}{1F4E79}
\definecolor{coulProp}{HTML}{7030A0}
\definecolor{coulRem}{HTML}{C55A11}
\definecolor{coulEx}{HTML}{2E7D32}
\definecolor{coulExo}{HTML}{B03A2E}
\definecolor{coulPart}{HTML}{1F4E79}
\definecolor{coulMeth}{HTML}{1B6E4F}
\definecolor{coulCourbe}{HTML}{0B5ED7}
\definecolor{coulCourbeB}{HTML}{C00000}
\definecolor{coulCourbeC}{HTML}{2E7D32}
\definecolor{coulGrille}{HTML}{8C8C8C}
\definecolor{coulRep}{HTML}{1B6E4F}
\definecolor{coulBar}{HTML}{2E7D32}

\newcommand{\R}{\mathbb{R}}
\newcommand{\Rs}{\mathbb{R}^{*}}
\renewcommand{\le}{\leqslant}
\renewcommand{\ge}{\geqslant}
\newcommand{\vect}[1]{\overrightarrow{#1}}
\newcommand{\norme}[1]{\left\lVert #1 \right\rVert}
\newcommand{\intervalle}[2]{\left[#1\,;#2\right]}
\RequirePackage{cancel}
\renewcommand{\CancelColor}{\color{coulExo}}
\newcommand{\fois}[1]{\times{\color{coulExo}#1}}
\newcommand{\barre}[1]{\cancel{#1}}

\tikzset{grille/.style={coulGrille,line width=0.4pt}}
\newcommand{\quadrillage}[4]{\draw[grille,step=1] (#1,#2) grid (#3,#4);}
\tikzset{
  croix/.style={cross out,draw,line width=0.6pt,inner sep=0pt,minimum size=4pt},
  croixdroite/.style={croix,rotate=45,line width=0.8pt},
}
\newcommand{\pt}[3]{\path (#1) node[croix]{} node[#2,font=\small,inner sep=2.5pt]{$#3$};}

\newcommand{\lycee@repere}[4]{%
  \draw[grille,step=1] (#1,#2) grid (#3,#4);
  \draw[-{Stealth[length=2mm]},thick] (#1,0)--({#3+0.4},0) node[below left=-1pt]{$x$};
  \draw[-{Stealth[length=2mm]},thick] (0,#2)--(0,{#4+0.4}) node[below left=-1pt]{$y$};
  \node[below left,font=\scriptsize,inner sep=1.5pt] at (0,0) {$0$};}
\newcommand{\lycee@gradx}[1]{\foreach \k/\l in {#1}{%
  \draw (\k,0.07)--(\k,-0.07) node[below,font=\scriptsize,inner sep=1.5pt]{$\l$};}}
\newcommand{\lycee@grady}[1]{\foreach \k/\l in {#1}{%
  \draw (0.07,\k)--(-0.07,\k) node[left,font=\scriptsize,inner sep=1.5pt]{$\l$};}}
\newcommand{\lycee@intff}[2]{\left[#1\,;#2\right]}
\newcommand{\lycee@intfo}[2]{\left[#1\,;#2\right[}
\newcommand{\lycee@intof}[2]{\left]#1\,;#2\right]}
\newcommand{\lycee@intoo}[2]{\left]#1\,;#2\right[}
\newcommand{\lycee@ens}[1]{\left\{#1\right\}}
\AtBeginDocument{%
  \providecommand{\repere}{\lycee@repere}%
  \providecommand{\gradx}{\lycee@gradx}%
  \providecommand{\grady}{\lycee@grady}%
  \providecommand{\Cf}{\mathcal{C}\sb{\scriptscriptstyle f}}%
  \providecommand{\Cg}{\mathcal{C}\sb{\scriptscriptstyle g}}%
  \providecommand{\intff}{\lycee@intff}%
  \providecommand{\intfo}{\lycee@intfo}%
  \providecommand{\intof}{\lycee@intof}%
  \providecommand{\intoo}{\lycee@intoo}%
  \providecommand{\ens}{\lycee@ens}}

\newlist{questions}{enumerate}{3}
\setlist[questions,1]{label=\arabic*\textdegree),leftmargin=*,itemsep=2pt,topsep=3pt}
\setlist[questions,2]{label=\alph*),leftmargin=*,itemsep=1pt,topsep=2pt}
\setlist[questions,3]{label=\roman*),leftmargin=*,itemsep=1pt,topsep=1pt}
\setlist[itemize]{label=\textbullet,leftmargin=*,itemsep=2pt,topsep=3pt}

\newcommand{\besoinplace}[1]{\par\penalty\@M\begingroup
  \dimen@=#1\relax\dimen@ii=\pagegoal\advance\dimen@ii-\pagetotal
  \ifdim\dimen@>\dimen@ii\ifdim\dimen@ii>\z@\vfil\fi\break\fi\endgroup}

\newcommand{\lycee@pied}{\thepage/\pageref{LastPage}}
\newcommand{\entete}[2]{%
  \pagestyle{fancy}\fancyhf{}%
  \fancyhead[L]{\footnotesize\color{coulPart}\textbf{#1}}%
  \fancyhead[R]{\footnotesize\color{coulPart}\textbf{#2}}%
  \fancyfoot[C]{\footnotesize\lycee@pied}%
  \renewcommand{\headrulewidth}{0.4pt}%
  \renewcommand{\headrule}{\hbox to\headwidth{%
    \color{coulPart!40}\leaders\hrule height \headrulewidth\hfill}}}

\RequirePackage{ccicons}
\newcommand{\lycee@licence}{{\scriptsize\color{gray}%
  \copyright\ Luc Destombe --- \ccbyncsa\ CC BY-NC-SA 4.0}}
\newif\iflycee@licence \lycee@licencetrue
\newcommand{\sanslicence}{\lycee@licencefalse}
\AtBeginDocument{\iflycee@licence\AddToHookNext{shipout/foreground}{%
  \put(0,0){\rlap{\kern\dimexpr1in+\hoffset+\oddsidemargin+\textwidth\relax
    \llap{\raisebox{-\dimexpr1in+\voffset+\topmargin+\headheight+\headsep
      +\textheight+\footskip\relax}{\lycee@licence}}}}}\fi}

\newcommand{\formulesserrees}{%
  \setlength{\abovedisplayskip}{3pt}\setlength{\belowdisplayskip}{3pt}%
  \setlength{\abovedisplayshortskip}{2pt}\setlength{\belowdisplayshortskip}{3pt}}

\setlength{\parindent}{0pt}

\RequirePackage[implicit=false,hidelinks,pdfencoding=auto,psdextra,bookmarksopen,
  bookmarksopenlevel=1,pdfcreator={},pdfproducer={}]{hyperref}
\RequirePackage{bookmark}
\hypersetup{pdfauthor={Luc Destombe},pdfkeywords={CC BY-NC-SA 4.0}}
\newcounter{lycee@signet}
\newcommand{\signet}[3][]{\stepcounter{lycee@signet}%
  \edef\lycee@dest{\if\relax\detokenize{#1}\relax
    signet.\arabic{lycee@signet}\else#1\fi}%
  \hypertarget{\lycee@dest}{}\bookmark[level=#2,dest=\lycee@dest]{#3}}
\pdfstringdefDisableCommands{%
  \def\R{R}\def\vect#1{#1}\def\Cf{Cf}\def\Cg{Cg}\def\fois#1{×#1}%
  \def\barre#1{#1}\def\intervalle#1#2{[#1;#2]}\def\intff#1#2{[#1;#2]}%
  \def\textsuperscript#1{#1}\def\dfrac#1#2{#1/#2}\def\frac#1#2{#1/#2}%
  \def\vec#1{#1⃗}}%
\RequirePackage[official]{eurosym}

\geometry{top=0.8cm,bottom=1.3cm,left=1cm,right=1cm,footskip=0.6cm}

\pagestyle{fancy}
\fancyhf{}
\renewcommand{\headrulewidth}{0pt}
\fancyfoot[C]{\footnotesize Page \thepage{} / \pageref{LastPage}}

\newcommand{\titrefiche}[2]{%
  \begin{center}
    \Large\textbf{#1}\\[2pt]
    \large\textbf{#2}\\[2pt]
  \end{center}
  \vspace{0.10cm}}

\setlength{\columnseprule}{0.4pt}
\setlength{\columnsep}{15pt}
\renewcommand{\columnseprulecolor}{\color{black!45}}
\setlength{\multicolsep}{2pt}

\newcounter{partiefiche}
\renewcommand{\thepartiefiche}{\Roman{partiefiche}}
\newif\ifapresbandeau
\newcommand{\lycee@bandeau}[1]{%
  \par\penalty-600
  \vspace{0.08cm}%
  \noindent\colorbox{coulPart}{%
    \parbox{\dimexpr\linewidth-2\fboxsep\relax}{%
      \color{white}\bfseries\raggedright\strut#1\strut}}%
  \nopagebreak\par\nopagebreak
  \vspace{0.06cm}\nopagebreak
  \apresbandeautrue}
\newcommand{\partiefiche}[1]{%
  \refstepcounter{partiefiche}%
  \lycee@bandeau{\signet{1}{\thepartiefiche{} --- #1}\thepartiefiche{} --- #1}}
\newcommand{\partiefichesansnum}[1]{\lycee@bandeau{\signet{1}{#1}#1}}

\newcounter{exo}
\newlength{\espaceexo}\setlength{\espaceexo}{7pt}
\newlength{\espacetitre}\setlength{\espacetitre}{2pt}
\newcommand{\exo}[1][\relax]{%
  \par
  \ifapresbandeau\apresbandeaufalse\else\penalty-150\addvspace{\espaceexo}\fi
  \refstepcounter{exo}%
  \noindent\signet[exercice-\arabic{exo}]{2}{Exercice \arabic{exo}}%
  \textbf{\color{coulExo}\underline{Exercice \theexo}}%
  \ifx\relax#1\relax\else\textbf{ : #1}\fi
  \par\nopagebreak\vspace{\espacetitre}\nopagebreak
  \ignorespaces}

\setlist[enumerate]{nosep,leftmargin=*,itemsep=2pt,topsep=2pt}
\setlist[itemize]{nosep,leftmargin=*,topsep=2pt}
\newcommand{\choix}[3]{\par\hspace*{1em}%
  \makebox[0.3\linewidth][l]{a) #1}\makebox[0.3\linewidth][l]{b) #2}c) #3}
\newcommand{\deux}[2]{\par\hspace*{0.6em}%
  \makebox[0.48\linewidth][l]{#1}#2}
\newcommand{\trois}[3]{\par\hspace*{0.6em}%
  \makebox[0.32\linewidth][l]{#1}\makebox[0.32\linewidth][l]{#2}#3}

  \renewcommand{\exo}[1][\relax]{%
    \ifapresbandeau\apresbandeaufalse\else\par\penalty-150\fi
    \refstepcounter{exo}%
    \vspace{0.05cm}%
    \noindent\signet[exercice-\arabic{exo}]{2}{Exercice \arabic{exo}}%
    \textbf{\color{coulExo}\underline{Exercice \theexo}}%
    \ifx\relax#1\relax\else\textbf{ : #1}\fi%
    \par\nopagebreak
    \ignorespaces}
  \newcommand{\rappel}[1]{%
    \par\vspace{0.06cm}\nopagebreak
    \noindent\fcolorbox{black!35}{black!5}{%
      \parbox{\dimexpr\linewidth-2\fboxsep-2\fboxrule\relax}{%
        \small\textbf{Rappel.} #1}}%
    \par\vspace{0.06cm}\nopagebreak}
  \setlist[enumerate]{nosep,leftmargin=*}
  \setlist[itemize]{nosep,leftmargin=*}
  \setlength{\multicolsep}{3pt plus 1pt minus 1pt}
  \setlength{\premulticols}{10pt}
  \setlength{\postmulticols}{10pt}
  \newlist{expr}{enumerate}{1}
  \setlist[expr]{label={\Alph*\ $=$},nosep,leftmargin=*,itemsep=0pt}

\renewenvironment{center}{\par\addvspace{3pt}\centering}{\par\addvspace{3pt}}
\formulesserrees
\AtBeginDocument{\formulesserrees}
\clubpenalty=10000
\widowpenalty=10000
\displaywidowpenalty=10000
\brokenpenalty=10000
\setlength{\parskip}{0pt}

\RequirePackage{ragged2e}
\setlength{\RaggedRightRightskip}{0pt plus 1fil}
\AtBeginDocument{\RaggedRight\binoppenalty=10000 \relpenalty=10000 }
\g@addto@macro\@parboxrestore{\RaggedRight}
\makeatother
\usepackage{tkz-tab}
\definecolor{coulCourbe}{HTML}{C0392B}
\definecolor{coulCourbeB}{HTML}{1F4E79}

\begin{document}

\titrefiche{1\textsuperscript{re} STI2D --- Fonctions}{Fiche d'exercices du chapitre}

\begin{multicols}{2}\raggedcolumns

\partiefiche{Antécédent et image}

\exo[Images et antécédent]
On considère la fonction $f$ définie par $f(x)=4x-2x^{2}$.
\begin{enumerate}[label=\arabic*)]
  \item Calculer l'image de $1$, de $-2$, puis de $2{,}5$.
  \item Le nombre $3$ est-il un antécédent de $-6$ par $f$ ? Justifier.
\end{enumerate}

\exo[Une fonction affine]
On considère la fonction $g$ définie par $g(x)=5x-8$.
\begin{enumerate}[label=\arabic*)]
  \item Calculer l'image de $-3$, de $4$, puis de $\dfrac{6}{5}$.
  \item Le nombre $\dfrac{2}{5}$ est-il un antécédent de $6$ par $g$ ? Justifier.
\end{enumerate}

\exo[Plusieurs antécédents]
On considère la fonction $h$ définie par $h(x)=x^{2}-x-6$.
\begin{enumerate}[label=\arabic*)]
  \item Calculer $h(0)$, $h(2)$ et $h(-2)$.
  \item Le nombre $3$ est-il un antécédent de $0$ par $h$ ? Justifier.
  \item Citer deux antécédents de $0$ par $h$.
\end{enumerate}

\exo[Calculer des antécédents]
Pour chaque fonction $f$, calculer les antécédents du nombre $y$.
\begin{multicols}{2}
\begin{enumerate}[label=\alph*)]
  \item $f(x)=x-4$ ; \ $y=3$
  \item $f(x)=x^{2}$ ; \ $y=36$
  \item $f(x)=\dfrac{1}{x}$ ; \ $y=4$
  \item $f(x)=x(x-5)$ ; \ $y=0$
\end{enumerate}
\end{multicols}

\partiefiche{Ensemble de définition}

\exo[Valeurs interdites]
On considère les fonctions
\[
  f(x)=\dfrac{1}{x} \qquad g(x)=\dfrac{x+2}{x-5} \qquad h(x)=\sqrt{x}
\]
et les nombres $-2$ ; $5$ ; $0$ ; $9$. Pour chaque fonction, un de ces quatre nombres
n'a pas d'image : lequel ? Justifier.

\exo[Calculer, si c'est possible]
On considère les fonctions $f(x)=\dfrac{8}{x-4}$ et $g(x)=\sqrt{x+2}$. Calculer, si
c'est possible :
\begin{multicols}{2}
\begin{enumerate}[label=\alph*)]
  \item $f(6)$ ; $f(-4)$ ; $f(4)$
  \item $g(7)$ ; $g(14)$ ; $g(-3)$
\end{enumerate}
\end{multicols}

\exo[Fractions]
Déterminer l'ensemble de définition de chaque fonction.
\begin{multicols}{2}
\begin{enumerate}[label=\alph*)]
  \item $f(x)=\dfrac{1}{x+3}$
  \item $g(x)=\dfrac{x-2}{x-9}$
  \item $h(x)=\dfrac{5}{3x-12}$
  \item $k(x)=\dfrac{2x}{4x+20}$
\end{enumerate}
\end{multicols}

\exo[Racines carrées]
Déterminer l'ensemble de définition de chaque fonction.
\begin{multicols}{2}
\begin{enumerate}[label=\alph*)]
  \item $f(x)=\sqrt{x-4}$
  \item $g(x)=\sqrt{x+6}$
  \item $h(x)=\sqrt{2x-10}$
  \item $k(x)=\sqrt{3x+5}$
\end{enumerate}
\end{multicols}

\partiefiche{Tableau de valeurs, courbe et lecture graphique}

\exo[Première lecture graphique]
Voici la courbe d'une fonction $f$ définie sur $\intff{-2}{4}$.
\begin{center}
\begin{tikzpicture}[x=0.4cm,y=0.4cm]
  \repere{-3}{-5}{5}{5}
  \gradx{-2,-1,1,2,3,4} \grady{-4,-3,-2,-1,1,2,3,4}
  \draw[coulCourbe,line width=1.1pt,domain=-2:4,samples=60] plot (\x,{-\x*\x+2*\x+3});
  \node[coulCourbe,font=\footnotesize] at (3.8,2.6) {$\Cf$};
\end{tikzpicture}
\end{center}
Compléter par lecture graphique.
\begin{enumerate}[label=\alph*)]
  \item $f(0)=\dots$ \quad et \quad $f(2)=\dots$
  \item L'image de $-1$ par $f$ est $\dots$
  \item $4$ a pour antécédent $\dots$
  \item Les antécédents de $0$ par $f$ sont $\dots$ et $\dots$
  \item Le nombre $5$ a-t-il un antécédent par $f$ ? Justifier.
\end{enumerate}

\exo[Lire un tableau de valeurs]
Voici un tableau de valeurs d'une fonction $f$.
\begin{center}\small
\renewcommand{\arraystretch}{1.25}
\begin{tabular}{|c|c|c|c|c|c|}
\hline
$x$ & $-2$ & $-1$ & $0$ & $1$ & $2$\\
\hline
$f(x)$ & $4$ & $2$ & $-1$ & $-3$ & $1$\\
\hline
\end{tabular}
\end{center}
\begin{enumerate}[label=\arabic*)]
  \item Donner l'image de $-1$, de $1$, puis de $-2$.
  \item Le nombre $1$ est-il un antécédent de $-3$ ?
  \item Donner un antécédent de $1$.
\end{enumerate}

\exo[Images et antécédents dans un tableau]
Voici un tableau de valeurs d'une fonction $g$.
\begin{center}\small
\renewcommand{\arraystretch}{1.25}
\begin{tabular}{|c|c|c|c|c|c|}
\hline
$x$ & $-3$ & $-1$ & $2$ & $4$ & $6$\\
\hline
$g(x)$ & $-4$ & $6$ & $0$ & $-1$ & $-4$\\
\hline
\end{tabular}
\end{center}
\begin{enumerate}[label=\arabic*)]
  \item Donner l'image de $4$, de $2$, puis de $-1$.
  \item Citer un antécédent de $6$, de $-1$, puis de $-4$.
  \item Pourquoi dit-on « \emph{un} antécédent » et non « \emph{l'}antécédent » ?
\end{enumerate}

\exo[Lecture graphique complète]
Voici la courbe d'une fonction $f$.
\begin{center}
\begin{tikzpicture}[x=0.4cm,y=0.4cm]
  \repere{-5}{-3}{6}{3}
  \gradx{-4,-3,-2,-1,1,2,3,4,5} \grady{-2,-1,1,2}
  \draw[coulCourbe,line width=1.1pt]
    (-4,-2) .. controls (-3,-1) and (-2,0) .. (-1,1)
    .. controls (-0.667,1.333) and (-0.333,2) .. (0,2)
    .. controls (0.333,2) and (0.667,1.333) .. (1,1)
    .. controls (1.333,0.667) and (1.667,0.333) .. (2,0)
    .. controls (2.333,-0.333) and (2.667,-1) .. (3,-1)
    .. controls (3.333,-1) and (3.667,-0.333) .. (4,0)
    .. controls (4.333,0.333) and (4.667,0.667) .. (5,1);
  \path[coulCourbe] (-4,-2) node[croix]{} (5,1) node[croix]{};
  \node[coulCourbe,font=\footnotesize] at (5.4,2) {$\Cf$};
\end{tikzpicture}
\end{center}
Déterminer, par lecture graphique :
\begin{multicols}{2}
\begin{enumerate}[label=\alph*)]
  \item l'ensemble de définition de $f$ ;
  \item l'image de $3$ ;
  \item l'image de $-1$ ;
  \item $f(1)$ et $f(-4)$ ;
  \item les antécédents de $2$ ;
  \item les antécédents de $1$ ;
  \item les antécédents de $0$ ;
  \item les antécédents de $-1$ ;
  \item les antécédents de $3$ ;
  \item le maximum et le minimum de $f$, et où ils sont atteints.
\end{enumerate}
\end{multicols}

\exo[Deux expressions, une seule fonction]
On considère les fonctions $f$ et $g$ définies par
\[
  f(x)=(x+4)(3-x)+x^{2}-5 \qquad\text{et}\qquad g(x)=7-x .
\]
\begin{enumerate}[label=\arabic*)]
  \item Calculer $f(2)$ et $g(2)$.
  \item Choisir un autre nombre et calculer son image par $f$, puis par $g$.
  \item Démontrer que tout nombre a la même image par $f$ et par $g$.
\end{enumerate}

\partiefiche{Tableau de variation d'une fonction}

\exo[De la courbe au tableau]
La fonction $f$ est définie sur $\intff{-3}{5}$ par sa courbe ci-dessous. Dresser son
tableau de variation, puis donner son maximum et son minimum.
\begin{center}
\begin{tikzpicture}[x=0.4cm,y=0.4cm]
  \repere{-4}{-3}{6}{5}
  \gradx{-3,-2,-1,1,2,3,4,5} \grady{-2,-1,1,2,3,4}
  \draw[coulCourbe,line width=1.1pt]
    (-3,3) .. controls (-2,0.5) and (0,-2) .. (1,-2)
           .. controls (2,-2) and (4,1.5) .. (5,4);
  \foreach \p in {(-3,3),(1,-2),(5,4)} \path[coulCourbe] \p node[croix]{};
  \node[coulCourbe,font=\footnotesize] at (5.4,2.6) {$\Cf$};
\end{tikzpicture}
\end{center}

\exo[Décrire puis résumer]
Voici la courbe d'une fonction $f$ définie sur $\intff{-3}{4}$.
\begin{center}
\begin{tikzpicture}[x=0.4cm,y=0.4cm]
  \repere{-4}{-4}{5}{3}
  \gradx{-3,-2,-1,1,2,3,4} \grady{-3,-2,-1,1,2}
  \draw[coulCourbe,line width=1.1pt]
    (-3,-1) .. controls (-2.33,0.4) and (-1.67,2) .. (-1,2)
            .. controls (0,2) and (1,-3) .. (2,-3)
            .. controls (2.67,-3) and (3.33,-1.6) .. (4,0);
  \foreach \p in {(-3,-1),(-1,2),(2,-3),(4,0)} \path[coulCourbe] \p node[croix]{};
  \node[coulCourbe,font=\footnotesize] at (4.4,1.2) {$\Cf$};
\end{tikzpicture}
\end{center}
\begin{enumerate}[label=\arabic*)]
  \item Décrire par des phrases les variations de $f$.
  \item Dresser le tableau de variation de $f$.
\end{enumerate}

\exo[Du tableau à la courbe]
Voici le tableau de variation d'une fonction $f$.
\begin{center}
\begin{tikzpicture}[scale=0.8]
  \tkzTabInit[lgt=1.4,espcl=2.2]{$x$/1,$f$/2}{$-4$,$0$,$3$}
  \tkzTabVar{-/$-2$,+/$3$,-/$-1$}
\end{tikzpicture}
\end{center}
\begin{enumerate}[label=\arabic*)]
  \item Quel est l'ensemble de définition de $f$ ?
  \item Tracer dans un repère une courbe qui peut représenter $f$.
\end{enumerate}

\exo[Lire un tableau de variation]
Voici le tableau de variation d'une fonction $g$.
\begin{center}
\begin{tikzpicture}[scale=0.8]
  \tkzTabInit[lgt=1.4,espcl=1.5]{$x$/1,$g$/2}{$-5$,$-2$,$1$,$4$}
  \tkzTabVar{+/$3$,-/$-2$,+/$4$,-/$0$}
\end{tikzpicture}
\end{center}
\begin{enumerate}[label=\arabic*)]
  \item Quel est l'ensemble de définition de $g$ ?
  \item Décrire par des phrases les variations de $g$.
  \item Donner le maximum et le minimum de $g$, et où ils sont atteints.
  \item Tracer dans un repère une courbe qui peut représenter $g$.
\end{enumerate}

\exo[Construire un tableau à partir d'un texte]
La fonction $f$ est définie sur $\intff{-4}{5}$. On sait que :
\begin{itemize}
  \item $f$ est croissante sur $\intff{-4}{1}$ et sur $\intff{3}{5}$ ;
  \item $f$ est décroissante sur $\intff{1}{3}$ ;
  \item $f(-4)=-3$ ; \ $f(1)=f(5)=2$ ; \ $f(3)=-1$.
\end{itemize}
\begin{enumerate}[label=\arabic*)]
  \item Dresser le tableau de variation de $f$.
  \item Tracer dans un repère une courbe qui peut représenter $f$.
\end{enumerate}

\exo[Résolution graphique]
Voici la courbe d'une fonction $f$ définie sur $\intff{0}{10}$.
\begin{center}
\begin{tikzpicture}[x=0.4cm,y=0.4cm]
  \repere{-1}{-3}{11}{4}
  \gradx{1,2,3,4,5,6,7,8,9,10} \grady{-2,-1,1,2,3}
  \draw[coulCourbe,line width=1.1pt]
    (0,2) .. controls (1,1) and (2,0) .. (3,-1)
    .. controls (3.333,-1.333) and (3.667,-2) .. (4,-2)
    .. controls (4.333,-2) and (4.667,-1.333) .. (5,-1)
    .. controls (6,0) and (7,1) .. (8,2)
    .. controls (8.333,2.333) and (8.667,3) .. (9,3)
    .. controls (9.333,3) and (9.667,2.333) .. (10,2);
  \path[coulCourbe] (0,2) node[croix]{} (10,2) node[croix]{};
  \node[coulCourbe,font=\footnotesize] at (6.4,1.9) {$\Cf$};
\end{tikzpicture}
\end{center}
\begin{enumerate}[label=\arabic*)]
  \item Résoudre graphiquement :
        \begin{multicols}{3}
        \begin{enumerate}[label=\alph*)]
          \item $f(x)=2$
          \item $f(x)=1$
          \item $f(x)=-2$
          \item $f(x)=4$
          \item $f(x)=0$
        \end{enumerate}
        \end{multicols}
  \item Combien l'équation $f(x)=-1$ a-t-elle de solutions ?
  \item Résoudre graphiquement, en donnant les solutions sous forme d'intervalles :
        \begin{multicols}{2}
        \begin{enumerate}[label=\alph*)]
          \item $f(x)>-2$
          \item $f(x)<1$
          \item $f(x)\ge 1$
          \item $f(x)\ge 0$
        \end{enumerate}
        \end{multicols}
\end{enumerate}

\partiefiche{Fonctions affines --- équation réduite}

\exo[Reconnaître $m$ et $p$]
Chaque fonction est affine : $f(x)=mx+p$. Donner $m$ et $p$ pour chacune.
\begin{multicols}{2}
\begin{enumerate}[label=\alph*)]
  \item $f(x)=5x-2$
  \item $f(x)=-3x+6$
  \item $f(x)=x-4$
  \item $f(x)=7-3x$
  \item $f(x)=4$
  \item $f(x)=-2x$
  \item $f(x)=\dfrac{-x}{3}+1$
  \item $f(x)=\dfrac{2}{5}x-\dfrac{3}{5}$
\end{enumerate}
\end{multicols}

\exo[Images et ordonnée à l'origine]
On considère la fonction affine $f$ définie par $f(x)=-2x+5$.
\begin{enumerate}[label=\arabic*)]
  \item Calculer $f(4)$, puis $f(-3)$.
  \item Quelle est l'ordonnée à l'origine de la droite qui représente $f$ ?
\end{enumerate}

\exo[Pente et équation réduite (1)]
Déterminer la pente de la droite $(AB)$, puis son équation réduite, avec $A(5\,;-2)$ et
$B(0\,;3)$.

\exo[Pente et équation réduite (2)]
Même consigne avec $A(-2\,;1)$ et $B(4\,;10)$.

\partiefiche{Tracer la droite d'équation $y=mx+p$}

\exo[Sens de variation]
Donner, en justifiant, le sens de variation de chaque fonction affine.
\begin{multicols}{2}
\begin{enumerate}[label=\arabic*)]
  \item $f(x)=3x-1$
  \item $f(x)=-2x+7$
  \item $f(x)=x+4$
  \item $f(x)=5-x$
  \item $f(x)=\sqrt{2}\,(x-3)$
  \item $f(x)=\dfrac{4-3x}{5}$
\end{enumerate}
\end{multicols}

\exo[Associer chaque fonction à sa droite]
Associer chaque fonction à sa droite, en justifiant par la pente et l'ordonnée à
l'origine.
\begin{center}
\begin{tikzpicture}[x=0.32cm,y=0.32cm]
  \begin{scope}
    \repere{-3}{-4}{3}{4}
    \draw[coulCourbeB,line width=1.1pt,domain=-0.5:3] plot (\x,{2*\x-3});
    \node[font=\scriptsize] at (0,-5.3) {\textcircled{1}};
  \end{scope}
  \begin{scope}[xshift=2.5cm]
    \repere{-3}{-4}{3}{4}
    \draw[coulCourbeB,line width=1.1pt,domain=-3:3] plot (\x,{-\x/2+1});
    \node[font=\scriptsize] at (0,-5.3) {\textcircled{2}};
  \end{scope}
  \begin{scope}[xshift=5cm]
    \repere{-3}{-4}{3}{4}
    \draw[coulCourbeB,line width=1.1pt,domain=-2/3:2] plot (\x,{-3*\x+2});
    \node[font=\scriptsize] at (0,-5.3) {\textcircled{3}};
  \end{scope}
\end{tikzpicture}
\end{center}
\textit{Le quadrillage est unitaire.}
\[
  f(x)=-3x+2 \qquad g(x)=2x-3 \qquad h(x)=-\dfrac{1}{2}x+1
\]

\exo[Appartenance et tracé]
Soit $d$ la droite d'équation réduite $y=4x-3$.
\begin{enumerate}[label=\arabic*)]
  \item Vérifier que le point $A(1\,;1)$ est sur $d$.
  \item Trouver l'abscisse du point $B$ de $d$ dont l'ordonnée est $5$.
  \item Tracer la droite $d$.
\end{enumerate}

\exo[Tracer à partir d'un point et de la pente]
Pour chaque cas, tracer la droite qui passe par $A$ et de pente $m$, puis donner son
équation réduite.
\begin{multicols}{2}
\begin{enumerate}[label=\alph*)]
  \item $A(1\,;2)$ et $m=3$
  \item $A(3\,;-1)$ et $m=-2$
\end{enumerate}
\end{multicols}

\partiefiche{Équation cartésienne d'une droite}

\exo[Équation cartésienne de $(AB)$ (1)]
Déterminer une équation cartésienne de la droite $(AB)$, avec $A(2\,;3)$ et $B(5\,;1)$.
Contrôler le résultat avec les deux points.

\exo[Équation cartésienne de $(AB)$ (2)]
Même consigne avec $A(-3\,;-1)$ et $B(5\,;3)$.

\exo[Lire une équation cartésienne]
On considère la droite $d$ d'équation cartésienne $3x-4y+5=0$.
\begin{enumerate}[label=\arabic*)]
  \item Les points $A(1\,;2)$, $B(-3\,;-1)$ et $C(2\,;3)$ sont-ils sur $d$ ? Justifier
        par un calcul.
  \item Trouver l'ordonnée du point de $d$ d'abscisse $5$.
  \item Trouver l'abscisse du point de $d$ d'ordonnée $-4$.
  \item Donner un vecteur directeur de $d$.
\end{enumerate}

\partiefiche{Tracer une droite d'équation $ax+by+c=0$}

\exo[Point et vecteur directeur]
\begin{enumerate}[label=\arabic*)]
  \item Déterminer une équation cartésienne de la droite $d$ qui passe par $A(1\,;-2)$
        et de vecteur directeur $\vec{u}\,(1\,;3)$, puis tracer $d$.
  \item On donne $B(-4\,;-3)$ et $C(2\,;0)$. Déterminer une équation cartésienne de la
        droite $(BC)$, puis la tracer.
\end{enumerate}

\exo[Deux droites dans un même repère]
On considère les droites $d$ et $d'$ d'équations
\[
  2x+3y-6=0 \qquad\text{et}\qquad 4x-3y-3=0 .
\]
Tracer ces deux droites dans un même repère.

\textit{Chercher deux points de chaque droite, en choisissant des valeurs de $x$ qui
donnent un $y$ entier.}

\exo[Droites verticales et horizontales]
Tracer dans un même repère les droites d'équations :
\begin{multicols}{2}
\begin{enumerate}[label=\alph*)]
  \item $x+2=0$
  \item $y-4=0$
  \item $x-y-1=0$
  \item $3x+2y-4=0$
\end{enumerate}
\end{multicols}
Lesquelles sont verticales ? horizontales ? Lesquelles ont une équation réduite ?

\partiefiche{Passer d'une forme à l'autre}

\exo[Pente, équation réduite, équation cartésienne (1)]
On donne $A(-1\,;-4)$ et $B(2\,;2)$.
\begin{enumerate}[label=\arabic*)]
  \item Déterminer la pente de la droite $(AB)$, puis son équation réduite.
  \item En déduire une équation cartésienne de $(AB)$.
\end{enumerate}

\exo[Pente, équation réduite, équation cartésienne (2)]
Même consigne avec $A(0\,;3)$ et $B(4\,;1)$.

\exo[Deux droites à déterminer]
On donne les points $C(3\,;-2)$, $D(1\,;6)$ et $E(-1\,;5)$. Déterminer l'équation réduite :
\begin{enumerate}[label=\alph*)]
  \item de la droite $(CD)$ ;
  \item de la droite $d$ de pente $5$ qui passe par $E$.
\end{enumerate}

\exo[Un point et une pente]
Déterminer l'équation réduite de chaque droite, puis une équation cartésienne.
\begin{enumerate}[label=\alph*)]
  \item $d$ passe par $A(2\,;-3)$ et a pour pente $4$ ;
  \item $d'$ passe par $B(5\,;-4)$ et a pour pente $0$.
\end{enumerate}

\partiefiche{Fonction carré}

\exo[Images et antécédents par la fonction carré]
\begin{enumerate}[label=\arabic*)]
  \item Calculer le carré de $-4$, de $7$, de $-0{,}2$, puis de $20$.
  \item Donner les antécédents par la fonction carré de $25$, de $0$, de $-9$, puis
        de $11$.
\end{enumerate}

\exo[Se ramener à $x^{2}=k$]
\begin{enumerate}[label=\arabic*)]
  \item Montrer que l'équation $5x^{2}-4=-x^{2}+14$ équivaut à $x^{2}=3$.
  \item Résoudre dans $\R$ l'équation $x^{2}=3$.
  \item En déduire les solutions de $5x^{2}-4=-x^{2}+14$.
\end{enumerate}

\exo[Équations avec des carrés]
Résoudre dans $\R$ :
\begin{multicols}{2}
\begin{enumerate}[label=\alph*)]
  \item $3x^{2}-8=x^{2}+10$
  \item $6x^{2}-1=8x^{2}-11$
  \item $2x^{2}+7=3$
\end{enumerate}
\end{multicols}

\partiefiche{Fonction cube}

\exo[Images et antécédents par la fonction cube]
\begin{enumerate}[label=\arabic*)]
  \item Calculer le cube de $3$, de $-2$, de $0{,}2$, puis de $-5$.
  \item Donner l'antécédent par la fonction cube de $64$, de $-27$, de $1\,000$, puis
        de $6$.
\end{enumerate}

\exo[Équations avec des cubes (1)]
Résoudre dans $\R$ :
\begin{multicols}{2}
\begin{enumerate}[label=\alph*)]
  \item $-3x^{3}+10=5x^{3}-6$
  \item $5x^{3}-3=2x^{3}+21$
\end{enumerate}
\end{multicols}

\exo[Équations avec des cubes (2)]
Résoudre dans $\R$ :
\begin{multicols}{2}
\begin{enumerate}[label=\alph*)]
  \item $4x^{3}+9=x^{3}-15$
  \item $-x^{3}-4=2x^{3}+11$
\end{enumerate}
\end{multicols}

\partiefiche{Fonction inverse}

\exo[Lecture graphique sur l'hyperbole]
Voici la courbe de la fonction inverse et les droites d'équations $y=1$ et
$y=-\dfrac{1}{2}$.
\begin{center}
\begin{tikzpicture}[x=0.36cm,y=0.36cm]
  \repere{-5}{-4}{5}{4}
  \gradx{-4,-3,-2,-1,1,2,3,4} \grady{-3,-2,-1,1,2,3}
  \draw[coulCourbeB,thick] (-5,1) -- (5,1) node[above left,font=\scriptsize] {$y=1$};
  \draw[coulCourbeB,thick] (-5,-0.5) -- (5,-0.5)
    node[below left,font=\scriptsize] {$y=-\frac{1}{2}$};
  \draw[coulCourbe,line width=1.1pt,domain=0.25:5,samples=60] plot (\x,{1/\x});
  \draw[coulCourbe,line width=1.1pt,domain=-5:-0.25,samples=60] plot (\x,{1/\x});
\end{tikzpicture}
\end{center}
\begin{enumerate}[label=\arabic*)]
  \item Résoudre graphiquement $\dfrac{1}{x}=1$, puis $\dfrac{1}{x}=-\dfrac{1}{2}$.
  \item Résoudre graphiquement $\dfrac{1}{x}<1$, puis $\dfrac{1}{x}>-\dfrac{1}{2}$.
\end{enumerate}
\textit{$0$ est une valeur interdite.}

\exo[Équations et inéquations (1)]
Résoudre dans $\R$ :
\begin{multicols}{3}
\begin{enumerate}[label=\alph*)]
  \item $\dfrac{1}{x}=5$
  \item $\dfrac{1}{x}\le 5$
  \item $\dfrac{1}{x}>5$
\end{enumerate}
\end{multicols}

\exo[Équations et inéquations (2)]
Résoudre dans $\R$ :
\begin{multicols}{3}
\begin{enumerate}[label=\alph*)]
  \item $\dfrac{1}{x}=-2$
  \item $\dfrac{1}{x}\le -2$
  \item $\dfrac{1}{x}<-2$
\end{enumerate}
\end{multicols}

\exo[Comparer deux inverses]
Résoudre dans $\R$ :
\begin{multicols}{3}
\begin{enumerate}[label=\alph*)]
  \item $\dfrac{1}{x}\le\dfrac{1}{4}$
  \item $\dfrac{1}{x}\ge\dfrac{3}{5}$
  \item $\dfrac{1}{x}>\dfrac{1}{2}$
\end{enumerate}
\end{multicols}

\exo[Quotients du premier degré]
Résoudre chaque équation, après avoir donné la valeur interdite.
\begin{multicols}{2}
\begin{enumerate}[label=\alph*)]
  \item $\dfrac{6}{3x-3}=2$
  \item $\dfrac{-4}{-2x+8}=2$
\end{enumerate}
\end{multicols}

\partiefiche{Fonction racine carrée}

\exo[Racine carrée (1)]
Résoudre :
\begin{multicols}{3}
\begin{enumerate}[label=\alph*)]
  \item $\sqrt{x}=8$
  \item $\sqrt{x}\le 8$
  \item $\sqrt{x}>8$
\end{enumerate}
\end{multicols}

\exo[Racine carrée (2)]
Résoudre :
\begin{multicols}{3}
\begin{enumerate}[label=\alph*)]
  \item $\sqrt{x}=10$
  \item $\sqrt{x}\le 10$
  \item $\sqrt{x}<10$
\end{enumerate}
\end{multicols}
\textit{Comparer les réponses de b) et de c).}

\exo[Racine carrée (3)]
Résoudre :
\begin{multicols}{3}
\begin{enumerate}[label=\alph*)]
  \item $\sqrt{x}\le 7$
  \item $\sqrt{x}\ge 6$
  \item $\sqrt{x}>1$
\end{enumerate}
\end{multicols}

\exo[Inverses de carrés, de cubes et de racines]
Résoudre chaque équation, après avoir donné les valeurs interdites.
\begin{multicols}{3}
\begin{enumerate}[label=\alph*)]
  \item $\dfrac{4}{x^{2}}=1$
  \item $\dfrac{-2}{x^{3}}=16$
  \item $\dfrac{21}{\sqrt{x}}=3$
\end{enumerate}
\end{multicols}

\end{multicols}

\end{document}
